\documentclass[oneside,12pt]{book}

% Package dependencies
\usepackage{geometry}
\usepackage{lipsum}
\geometry{left=32mm, right=30mm, bottom=25mm, top=25mm}
\usepackage{amsmath , amsthm , amssymb}
\usepackage{float}
\usepackage{graphicx}
\usepackage{hyperref}
\usepackage{apacite}
\usepackage{scrextend}
\usepackage{blindtext}
\usepackage{fancyhdr}
\usepackage{multicol}
\usepackage{csquotes}
\usepackage{bm}
\usepackage{subcaption}
\usepackage{tocbibind}
\usepackage{pdfpages}
\usepackage[toc,page]{appendix}
\usepackage{wrapfig}
\usepackage{colortbl}
\usepackage{multirow}
\usepackage{soul}

% document formatting
\addtokomafont{labelinglabel}{\sffamily}
\setlength{\columnsep}{1cm}
\renewcommand{\baselinestretch}{1.5}


\begin{document}

% --- README ---
% The \hl command is used to highlight parts that need to be
% removed or replaced by your own content. Please read these
% carefully and make sure they are all correctly addressed.

% Title page
\begin{titlepage}
    \begin{center}
        \vspace*{1.5cm}
        
        \Huge
        \textbf{\hl{Thesis Title}}
        
        \vspace{0.5cm}
        \begin{figure}[H]
    	\centering
    	\hspace{7mm} \includegraphics[scale=0.5]{TU_logo}
        \end{figure}
        
        \vspace{1.5cm}
        
        \textbf{\hl{Student Name}}
        
        \vfill
       \large
        A dissertation submitted in partial fulfilment of the requirements of\\
	Technological University Dublin for the degree of\\
	M.Sc. in Computing \hl{(Stream)}\\
       \vspace{0.5cm}
        \textbf{\hl{Date}}
        \vspace{0.8cm}
 
    \end{center}
\end{titlepage}

% Define headers / footer style for the whole doc
\pagestyle{fancy}
\fancyhf{}
\fancyfoot[CE,CO]{\thepage}
\renewcommand{\headrulewidth}{0pt}

% APA style for referencing
\bibliographystyle{apacite}

% Roman numerals for the 'administrative' sections
\pagenumbering{Roman}

% Declaration page
\chapter*{Declaration}
\addcontentsline{toc}{chapter}{Declaration}
I certify that this dissertation which I now submit for examination for the award of
MSc in Computing \hl{(Stream)}, is entirely my own work and has not been taken
from the work of others save and to the extent that such work has been cited and
acknowledged within the text of my work.
\\
\\
This dissertation was prepared according to the regulations for postgraduate study of
the Technological University Dublin and has not been submitted in whole or part for an
award in any other Institute or University.
\\
\\
The work reported on in this dissertation conforms to the principles and requirements
of the Institute’s guidelines for ethics in research.
\vfill
\noindent
\textit{\textbf{Signed:}}  \\

\noindent
\textit{\textbf{Date:}}
\vspace{0.8cm}

% Abstract
\chapter*{Abstract}
\addcontentsline{toc}{chapter}{Abstract}
% --- Please remove the \hl command when writing your own abstract
\hl{Abstract outlining the project (max 750 words)
\\
\noindent
\textbf{Keywords:} \quad Keyword 1, keyword 2, ... (up to 8 keywords)}

\vfill

%Acknowledgements
\chapter*{Acknowledgments}
\addcontentsline{toc}{chapter}{Acknowledgments}
\hl{I would like to express my sincere thanks to...
(thank all the people how have assisted you in completing your dissertation. Start with your supervisor, all TU Dublin staff that may have helped, other people can include family and friends, industrial and academic staff from other institution, etc.)}


% Contents and lists
\newpage
\tableofcontents

\listoffigures

\listoftables

\chapter*{List of Acronyms}
\addcontentsline{toc}{chapter}{List of Acronyms}
\begin{table}[H]
  \centering
    \begin{tabular}{ l l }
    
    \hl{\textbf{ACRONYM1}} & \hl{Meaning} \\
    \hl{\textbf{ACRONYM2}} & \hl{Meaning} \\ 
    \end{tabular}
\end{table}

\newpage

% Standard numbering starts from here
\pagenumbering{arabic}
\fancyhead[RE,LO]{\leftmark}
\renewcommand{\headrulewidth}{2pt}


% Start of core thesis content

%Introduction chapter
% ============================================================= %
\chapter{Introduction}
% --- Please remove the content before the Background section --- %
\hl{\textbf{This template works only as a guideline to your dissertation. While it is highly recommended that you follow this suggested structure, it might not be possible in every case. Significant deviations will likely need to be discussed with your supervisor.}}
\\

\noindent \hl{This chapter should be a short account of why you undertook the
investigation, what the general state of knowledge was at the
time you started. Why you asked the questions that your
research/observations were expected to answer. It should state
your research question and briefly introduce the research
undertaken. A brief reader's guide to the dissertation should be 
included.}
\\

\noindent \hl{\textbf{Hints}}

\begin{itemize}
    \item \hl{In the introduction of each chapter cover the purpose of the chapter and give an overview of what it covers (expect for the introduction chapter)}
    \item \hl{At the end of each chapter summarise the main gist of the chapter itself and link it to the next chapter}
    \item \hl{Have a logical presentation of research, developments, findings and lead the reader in an integrated way to what you want to achieve}
    \item \hl{You can only claim to have an opinion if it is based on original research or an in-depth analysis, for instance, has been made of theory where different authors' viewpoints are contrasted (commonalities and differences), and you make deductions based on that.}
\end{itemize}

\section{Background}
% --- Please remove all the content in this section --- %
\par \hl{This content gives some examples of latex commands that might be used. To cite an article, using the apa style, at the end of a sentence, simply reference the label from the bibliography file. Examples below:}
\begin{itemize} 
\item \hl{phd thesis:} \cite{Longo2014}.
\item \hl{article in a conference proceeding:} \cite{Longo2016}
\item \hl{journal article:} \cite{einstein}
\item \hl{chapter book:} \cite{Longo2011c}
\item \hl{book:} \cite{dirac}
\item \hl{websites: references to web sites should be included as footnotes and not included in the References/Bibliography (they are not scientific contributions).}
\end{itemize}

\hl{If you wanna cite something at the beginning of the sentence see below.}

\begin{itemize}
    \item \citeA{Longo2016} \hl{proposed a framework for human mental workload representation and assessment.}
\end{itemize}

\par \hl{An example of a numbered bullet point list. To number the list, replace itemize with enumerate}
\begin{enumerate}
\item \hl{This is the first item.}
\item \hl{Item number 2.}
\item \hl{etc.}
\end{enumerate}

\par \hl{An example of a simple table. This can be referenced in the text using the label given to table} \ref{tab:sample_table}. \hl{The same can be done for sections, equations or figures.}
\begin{table}[H]
  \centering
    \caption{\hl{A sample table}}
    \begin{tabular}{| c | c | c |}
    \hline
     & \textbf{Column 1} & \textbf{Column 2} \\ \hline
    \textbf{Row 1} & 0 & 1 \\ \hline
    \textbf{Row 2} & 2 & 3 \\ \hline
    \textbf{Row 3} & 4 & 5 \\ \hline
    \end{tabular}
  \label{tab:sample_table}
\end{table}

\par \hl{An example of a figure or image:}

\begin{figure}[H]
    \centering
    \includegraphics[scale=0.5]{TU_logo}
    \caption{\hl{Sample image}}
    \label{fig:logo}
\end{figure}

\noindent \hl{\textbf{Tables and figures guidelines:}}
\begin{itemize}
    \item \hl{Tables and figures cannot span 2 pages.}
    \item \hl{Each figure and table has to have a caption.}
    \item \hl{Pages cannot be half empty because of a large figure or table cannot be placed in it.}
    \item \hl{Place the large figure or table in a new page and refer it with a label within the text (eg. Figure x.y in page x depicts...). }
    \item \hl{Each figure and table must be within left and right margins and must be fully readable}
\end{itemize}

\par \hl{A sample equation}:

\begin{equation}
e^{i\pi} = -1
\end{equation}

% -----------

\section{Research Project/problem}

\section{Research Objectives}

\section{Research Methodologies}

\section{Scope and Limitations}

\section{Document Outline}


% Literature review chapter
% ============================================================= %
\chapter{Review of existing literature}
% --- Please remove the paragraph below --- %
\hl{It is essential that this chapter should be a critical review in which the various papers are compared and in which you express your own opinion of the conclusions that may be drawn and to do your best to reconcile discrepant results in favour of one or other set. Provide a summary at the end of the sections or of the whole review. Remember that the content of this chapter must be relevant to the actual research carried out; it is not a ``brain dump" of everything you have read. You must demonstrate analysis and synthesis of the literature. Also in some cases it may be necessary to divide the state-of-the-art into two separate chapters; one covering the application domain, and the other the technologies, or one describing the background/context of your research and one on the state-of-art for your specific issue.}

\section{A section}

\subsection{A sub-section}

\subsubsection{A sub-sub-section}

\par Some content.


% Design and methodology chapter
% ============================================================= %
\chapter{Experiment design and methodology}
% --- Please remove the paragraph below --- %
\hl{The general structure of the study should be described clearly. The comparisons that are going to be made, the controls and technical details etc. should be included if appropriate. This chapter might report on the design of a software-based solution or an experiment using existing datasets. Ethical considerations should be described after the research design and before the data selection or collection procedures. Also the chapter includes the methodology/ies adopted for designing the solution and for evaluating it (e.g. errors, performance measures, accuracy, ROC curves, t-tests, correlations etc.)}

\section{A section}

\section {Ethical considerations}
% --- Please remove the content in this section --- %


\hl{You should try to demonstrate as much as possible your commitment to conducting research in an ethical and responsible manner. \textbf{If there are no critical issues identified, this should also be highlighted at this section.} Please refer to the MSc dissertation website for more information.}

\section{Another section}

\subsection{Another sub-section}

\subsubsection{Another sub-sub-section}

Some content

% Results, evaluation and discussion chapter
% ============================================================= %
\chapter{Results, evaluation and discussion}
% --- Please remove the paragraph below --- %
\hl{Depending on the nature of the project, this chapter will describe the actual work carried out e.g. any experiment undertaken or system implementation following the theoretical description of the design of Chapter 3 (design). Also it should focus on the presentation and discussion of the findings in the light of what is already supposed to be known, from the literature conducted in Chapter 2. It should show how findings confirm or refute the research hypothesis and how they differ with previous work in the literature. Do not use this section for another review of the literature.}

% Concluding chapter
% ============================================================= %
\chapter{Conclusion}
% --- Please remove the paragraph below --- %
\hl{This should be a short account of the results of your work, emphasising mainly what is new. There should be a close correlation between this chapter and Chapter 1, in which you described the problem you were addressing. It is advisable to deal with the limitations of your research at this stage and to suggest here what further work might be done. This is the appropriate place to do a self-assessment of your research.}

\section{Research Overview}

\section{Problem Definition}

\section{Design/Experimentation, Evaluation \& Results}

\section{Contributions and impact}

\section{Future Work \& recommendations}


% End of thesis content
% ============================================================= %

% Include the bibligraphy by referencing the correct .bib file
\bibliography{sample_bibliography}

% Optional appendices
\appendix
\chapter{Additional content}
% --- Please remove the paragraph below --- %
\hl{These should contain supplementary material that is not necessary in order for the reader to follow the argument. For example, the text of a questionnaire, detailed UML diagrams, or a complete Software Requirement Specification should be placed in an Appendix. It is not considered necessary to include code, but you may do so by including a link within your dissertation PDF.}

\end{document}